% \begingroup
%   \centering
%   \color{JHUBlue}

%   \vspace{-1cm}
%   \includegraphics[width=0.5\textwidth]{jhulogo-cropped.pdf}\par
  
%   \vspace*{.5cm}
%   {\Huge\textbf{Egbert Rijke}\par}
%   \vspace*{.3cm}
  
%   \begin{tikzpicture}[scale=4]
%     \draw[gray,ultra thin] (-1.1,-1.1) grid[step=0.2] (1.1,1.1);
%     \draw[ultra thick] (0,0) circle (1);
%     \coordinate (P1) at ({117/125},{44/125});
%     \coordinate (P2) at ({45/53},{-28/53});
%     \coordinate (P3) at ({5/13},{12/13});
%     \coordinate (P4) at ({-3/5},{4/5});
    
%     \draw[DarkOrchid,ultra thick,dashed] (1,1) -- (P1);
%     \draw[DarkOrchid,ultra thick,->,shorten >=2.4ex] (P1) -- (P2);
%     \draw[Green,ultra thick,dashed] (1,-1) -- (P2);
%     \draw[Green,ultra thick,->,shorten >=1ex] (P2) -- (P3);
%     \draw[YellowOrange,ultra thick,dashed] (1,1) -- (P3);
%     \draw[YellowOrange,ultra thick,->,shorten >=1.8ex] (P3) -- (P4);
%     \draw[CornflowerBlue,ultra thick,dashed] (-1,1) -- (P4);
%     \draw[CornflowerBlue,ultra thick,->,shorten >=1ex] (P4) -- (1,0);
%     \draw[fill=white,very thick] (P1) circle (.5pt);
%     \draw[fill=white,very thick] (P2) circle (.5pt);
%     \draw[fill=white,very thick] (P3) circle (.5pt); 
%     \draw[fill=white,very thick] (P4) circle (.5pt);
%     \draw[fill=JHUBlue,very thick] (1,0) circle (.5pt);
%   \end{tikzpicture}

%   \vspace*{.3cm}

%   {\superhuge \textbf{The Great Story of Numbers}\par}

%   \vspace*{.3cm}

%   {\LARGE Math 304\par}

%   \vspace*{.5cm}

%   {\Large\Version{0}\par}
% \endgroup

%%%% No Logo Version
%
\begin{titlepage}
  \begin{center}
    {\Huge\textsc{Johns Hopkins University}\par}
    \vspace*{\stretch{1}}
    {\Huge\textbf{\DocumentAuthor}\par}
    \vspace*{\stretch{1}}
    \begin{center}
    \begin{tikzpicture}[
      scale=5,
      alt={
        The unit circle on a square coordinate grid, centered at the origin.
        On the unit circle, a sequence of four Vieta jumps 
        is shown with arrows pointing to the target of each jump.
        This sequence of jumps ends at the point (1,0).
        Each jump is constructed using a dashed line segment
        from the nearest corner point (±1,±1) to the current point.
        Extending this line segment, we find a second intersection point on the
        unit circle, which is the target point of the jump.
        }
      ]
      \draw[gray,ultra thin] (-1.1,-1.1) grid[step=0.2] (1.1,1.1);
      \draw[ultra thick] (0,0) circle (1);
      \coordinate (P1) at ({117/125},{44/125});
      \coordinate (P2) at ({45/53},{-28/53});
      \coordinate (P3) at ({5/13},{12/13});
      \coordinate (P4) at ({-3/5},{4/5});

      \draw[DarkOrchid,ultra thick,dashed] (1,1) -- (P1);
      \draw[DarkOrchid,ultra thick,->,shorten >=2.4ex] (P1) -- (P2);
      \draw[Green,ultra thick,dashed] (1,-1) -- (P2);
      \draw[Green,ultra thick,->,shorten >=1ex] (P2) -- (P3);
      \draw[YellowOrange,ultra thick,dashed] (1,1) -- (P3);
      \draw[YellowOrange,ultra thick,->,shorten >=1.8ex] (P3) -- (P4);
      \draw[CornflowerBlue,ultra thick,dashed] (-1,1) -- (P4);
      \draw[CornflowerBlue,ultra thick,->,shorten >=1ex] (P4) -- (1,0);
      \draw[fill=white,very thick] (P1) circle (.5pt);
      \draw[fill=white,very thick] (P2) circle (.5pt);
      \draw[fill=white,very thick] (P3) circle (.5pt); 
      \draw[fill=white,very thick] (P4) circle (.5pt);
      \draw[fill=black,very thick] (1,0) circle (.5pt);
    \end{tikzpicture}
    \end{center}
    \vspace*{\stretch{.5}}
    \par\noindent
    \makebox[\textwidth][c]{%
      \superhuge \textbf{\SeriesTitle}%
    }
    \par
  \end{center}
\end{titlepage}
